\documentclass[a4paper,11pt,DIV=25]{scrartcl}

\usepackage[utf8x]{inputenc}
\usepackage[ngerman]{babel}
\usepackage{longtable}
%\usepackage{fancyhdr}
\usepackage{graphicx}
\usepackage{tabularx}
\usepackage{amsmath}
\usepackage{hyperref}
\author{Tim Reddehase}
\date{2014-09-09}

\newenvironment{oitemize}{%
  \renewcommand{\labelitemii}{$\bullet$}%
  \vspace{-15pt}
  \begin{itemize}}
  {\end{itemize}}

\begin{document}

\Huge
\begin{center}
  Ontohub loc/id\\
  \bigskip
  \includegraphics[width=0.3\textwidth]{ontohub-oor-logo.pdf}
\end{center}
\normalsize

\bigskip

\section{Idea}

Generally an ontohub loc/id (locator/identifier) is just an iri of an ontology
or one of its members (symbols, sentences, mappings, child-ontologies). However
ontohub loc/ids are generated by the ontohub application and assigned to an
ontology.  We somehow try to infer them from the path of the repository, the
path of the ontology and the specific name. Additionally we ensure that this
specific IRI is actually a locator and not \textit{just} an identifier.

\section{Status Report}

Currently ontologies have a somewhat MMT-like approach with a $?$ being used as
a separator between parent and child-ontology:

\begin{itemize}
  \item \url{http://ontohub.org/dol-testing/double_mapped_blendoid?DMB-CommonSource}
\end{itemize}

Currently this whole string of text is stored as a database field in the
ontologies-table.  This specifically includes the IRI-Schema, the hostname and
the implicit port.

\section{MMT-Style - Ontohub Edition}

Employing MMT in Ontohub would result in using $?$ as a separator (see the
Status Report).  However employing $?$ denotes the HTTP-Query-String. So
specifiying an ontology(-element) would use both the hierarchical part and part
of the query-string for IRIs.

\subsection{Examples}

\begin{tabularx}{\textwidth}{p{.1\textwidth}p{.2\textwidth}p{.7\textwidth}}
  \textit{Distributed Ontology} & & \\
  \hline
  & Distributed Ontology & \url{http://ontohub.org/dol-testing/double_mapped_blendoid}\\
  & Child Ontology & \url{http://ontohub.org/dol-testing/double_mapped_blendoid?DMB-CommonSource}\\
  & Mapping & \url{http://ontohub.org/dol-testing/double_mapped_blendoid?SomeMapping}\\
  & Symbol & \url{http://ontohub.org/dol-testing/double_mapped_blendoid?DMB-CommonSource?KitchenTable}\\
  & Sentence & \url{http://ontohub.org/dol-testing/double_mapped_blendoid?DMB-CommonSource?Ax02}\\
  & & \\
  \textit{Ontology} & & \\
  \hline
  & Distributed Ontology & \url{http://ontohub.org/dol-testing/double_mapped_blendoid}\\
  & Ontology & \url{http://ontohub.org/default/pizza}\\
  & Mapping & \url{http://ontohub.org/default/pizza?SomeMapping}\\
  & Symbol & \url{http://ontohub.org/default/pizza??KitchenTable}\\
  & Sentence & \url{http://ontohub.org/default/pizza??Ax02}\\
\end{tabularx}

\subsection{Conflicts in Names}

Technically the loc/id of a symbol and a sentence could conflict. Additionally
names of single ontologies in distributed ontologies (child-ontology) could
conflict with mappings.

\subsection{Benefits and Disadvantages}

The major benefit is of course that this system would be compatible with all
software products which comply with the (rather small) MMT-Standard.  Sadly
this breaks with some conventions across the web. For example data-retrieval in
different formats would solely rely on the MIME-Type of the HTTP requests
(Accept Header).

So the convention for Web-Application (and/or their APIs) to allow adding the
file-extension to a request in order to allow making file-based decision
without forcing specific Accept-Headers.

Technically we could still add this file-extension in the query-string, but
this would require special handling and still breaks the convention.

\section{Proposed Ontohub-Style}

\end{document}
